\section{Basic concepts}

1 Should be easy, but does it match what you expect?
2 So the spectra and the residual form a partition of C 
5 This already shows how functional analysis is going to be a lot harder. In finite dimensions, any operator/matrix always has a (possibly complex) eigenvalue, but here it is possible that there are no eigenvalues at all! On the other hand, the spectrum always is nonempty, so the spectrum takes over the role of the iegenvalues 
6
7
8 
10 Some useful inclusions of the spectrum

\begin{exercise}{1 (Identity operator)}
For the identity operator $I$ on a normed space $X$, find the eigenvalues and eigenspaces as well as $\sigma (I)$ and $R_\lambda(I)$.
\end{exercise}
\begin{proof}
fill
\end{proof} 

\begin{exercise}{2}
Show that for a given linear operator $T$, the sets $\rho(T), \sigma_p(T), \sigma_c(T)$, and $\sigma_r(T)$ are mutually disjoint and their union is the complex plane.
\end{exercise}
\begin{proof}
fill
\end{proof} 

\begin{exercise}{5}
Let $(e_k)$ be a total orthonormal sequence in a separable Hilbert space $H$ and let $T: H \to H$ be defined at $e_k$ by $T e_k = e_{k+1}$ for $k = 1, 2, \dots$, and then linearly and continously extend to $H$.
Find invariant subspaces.
Show that $T$ has no eigenvalues.
\end{exercise}
\begin{proof}
fill
\end{proof} 

\begin{exercise}{6 (Extension)}
The behaviour of the various parts of the spectrum under extension of an operator is of practical interest.
If $T$ is a bounded linear operator and $T'$ is a linear extension of $T$, show that we have $\sigma_p(T) \subseteq \sigma_p(T')$ and for any $\lambda \in \sigma_p(T)$ the eigenspace of $T$ is contained in the eigenspace of $T'$.
\end{exercise}
\begin{proof}
fill
\end{proof} 

\begin{exercise}{7}
Show that $\sigma_r(T') \subseteq \sigma_r(T)$ in exercise 6.
\end{exercise}
\begin{proof}
fill
\end{proof} 

\begin{exercise}{8}
Show that $\sigma_c(T) \subseteq \sigma_c(T') \cup \sigma_p(T')$ in exercise 6.
\end{exercise}
\begin{proof}
fill
\end{proof} 

\begin{exercise}{10}
How does $\rho(T') \subseteq \rho(T) \cup \sigma_r(T)$ follow from exercises 6 and 8?
\end{exercise}
\begin{proof}
fill
\end{proof} 
